\section{The integer program used as a cross-check}\label{sec:ip}


The leaf condition $\LEAF$ is only sufficient for the pairing condition $\CC(B,F,k)$ of \cref{def:C}.
To test the conclusion $\CC(B,F,k)$ directly, and independently of the max-flow machinery, we used a
$0$--$1$ program whose optimal value measures the largest violation of Hall's condition that Maker can
reach. It is solved in floating-point arithmetic, so its answers are computed statements of a weaker kind
than those of \cref{sec:verification}; \emph{none of them is used in the proof of \cref{thm:main}}.

\subsection{The program}\label{ssec:ip-model}

Let $B$ and $F$ be disjoint sets of cells and $k\ge 0$. Define
\[
\worst(B,F,k)=\max\bigl\{\,2|S|-|\N(S)\setminus X|\;:\; S\subseteq\live{B},\ X \text{ $B$-admissible},\
F\subseteq X,\ |X|\le k\,\bigr\},
\]
the maximum being taken over the empty family if no such $X$ exists. Since $S=\emptyset$ is allowed,
$\worst(B,F,k)\ge0$ whenever some $X$ exists, and by \cref{lem:hall} the pairing condition $\CC(B,F,k)$ holds
if and only if $\worst(B,F,k)\le 0$ (or no admissible $X$ exists, in which case $\CC$ holds vacuously).

The program has four families of binary variables: $s_l$ and $q_l$ for every live line $l\in\live{B}$, and
$n_c$ and $x_c$ for every cell $c\notin B$. The intended meaning is $s_l=1$ iff $l\in S$; $x_c=1$ iff
$c\in X$; $n_c=1$ for the cells of $\N(S)\cup X$; $q_l=1$ if $l$ carries four stones of $X$.
\begin{align}
\text{maximize}\quad & 2\sum_{l}s_l-\sum_{c}n_c+\sum_{c}x_c \label{eq:ip-obj}\\
\text{subject to}\quad
 & n_c\ge s_l && \text{for } l\in\live{B},\ c\in l, \label{eq:ip-n}\\
 & x_c\le n_c && \text{for } c\notin B, \label{eq:ip-xn}\\
 & \textstyle\sum_{c}x_c\le k,\qquad x_c=1 && \text{for } c\in F, \label{eq:ip-k}\\
 & \textstyle\sum_{c\in l}x_c\le 4 && \text{for } l\in\live{B}, \label{eq:ip-cap}\\
 & \textstyle q_l\ge\sum_{c\in l}x_c-3 && \text{for } l\in\live{B}, \label{eq:ip-q}\\
 & q_a+q_b-x_c\le 1 && \text{for } c\notin B \text{ and live lines } a\neq b \text{ through } c. \label{eq:ip-danger}
\end{align}

\begin{proposition}[Exactness of the integer program]\label{prop:ip-exact}
The program \eqref{eq:ip-obj}--\eqref{eq:ip-danger} is feasible if and only if $F$ is $B$-admissible and
$|F|\le k$, and then its optimal value is $\worst(B,F,k)$.
\end{proposition}

\begin{proof}
Let $(s,n,x,q)$ be feasible and put $X=\{c:x_c=1\}$, $S=\{l:s_l=1\}$. By \eqref{eq:ip-k}, $F\subseteq X$ and
$|X|\le k$; $X\cap B=\emptyset$ because there are no variables for the cells of $B$. By \eqref{eq:ip-cap}
every live line has at most four cells of $X$. If two live lines $a\ne b$ with exactly four cells of $X$
each share a cell $c$, then $q_a=q_b=1$ by \eqref{eq:ip-q}, hence $x_c=1$ by \eqref{eq:ip-danger}: the
common cell is in $X$, so $a,b$ is not a danger pair. Thus $X$ is $B$-admissible (\cref{def:admissible}).
Conversely, for a $B$-admissible $X$ with $F\subseteq X$, $|X|\le k$, and any $S\subseteq\live{B}$, set $x$ and
$s$ accordingly, $q_l=1$ exactly for the live lines with four cells of $X$, and $n_c=1$ exactly for
$c\in\N(S)\cup X$. All constraints hold; for \eqref{eq:ip-danger} note that $q_a=q_b=1$ forces the common
cell into $X$, because $X$ has no danger pair. So the feasible choices of $(x,s)$ are exactly the pairs
$(X,S)$ in the definition of $\worst$; in particular the program is feasible iff $X=F$ is feasible, i.e.\ iff
$F$ is $B$-admissible and $|F|\le k$.

For fixed $(X,S)$ the constraints \eqref{eq:ip-n} and \eqref{eq:ip-xn} say that the set $\{c:n_c=1\}$
contains $\N(S)\cup X$, and the objective is $2|S|-|\{c:n_c=1\}\setminus X|$, which is largest, namely
$2|S|-|\N(S)\setminus X|$, when equality holds. The variables $q$ do not occur in the objective. Hence the
optimal value is the maximum of $2|S|-|\N(S)\setminus X|$ over the admissible pairs $(X,S)$.
\end{proof}

The constraints \eqref{eq:ip-q} and \eqref{eq:ip-danger} encode the danger rule. Without them the
program computes the analogous maximum over all sets $X$ that merely respect the cap of four stones per
live line; the difference is visible already for $\Teight$ (\cref{tab:ip-regression}).

For $B=\Teight$ the program has $2\cdot100+2\cdot208=616$ binary variables and $1{,}605$ constraints, $596$
of which are of type \eqref{eq:ip-danger}.

\subsection{Solver and settings}\label{ssec:ip-solver}

The program is built by \fn{py/ip_check.py}, which has its own implementation of the geometry
(\fn{py/geom.py}) and shares no code with the C++ programs, and is solved with HiGHS~\cite{HuangfuHall2018},
version~1.16, through its Python interface, on one thread. The objective is integral, so the relative gap
is set to $0$ and the absolute gap to $0.99$: a gap below $1$ proves optimality of the incumbent. Two
modes are used.
\begin{itemize}[leftmargin=1.6em]
\item \emph{Full mode} computes $\worst(B,F,k)$.
\item \emph{Decide mode} only decides whether $\worst(B,F,k)\ge1$: the solver is told to stop as soon as it
  has a solution with objective value at least $1$ (objective target $0.999$); otherwise it runs until it
  has proved that the optimum is below $1$.
\end{itemize}
Before the solver is called, the program checks in exact arithmetic whether $F$ is $B$-admissible and
$|F|\le k$; if not, the instance is reported as infeasible (the pairing condition holds vacuously). Every
solution returned by the solver is re-evaluated in exact integer arithmetic: $X$ and $S$ are read off,
admissibility, $F\subseteq X$ and $|X|\le k$ are checked, and $2|S|-|\N(S)\setminus X|$ is recomputed from the
sets. Consequently an answer ``fails'' ($\worst\ge1$) comes with an integer-verified violating position and
is exact. An answer ``holds'' ($\worst\le0$) cannot be verified from a solution: it rests on the dual bound
reported by a floating-point branch-and-bound code. This asymmetry is the reason why the integer program
is used as a cross-check only.

\subsection{Regression values}\label{ssec:ip-regression}

\Cref{tab:ip-regression} lists values of $\worst$ that were fixed in advance as regression targets and
that the program reproduces (file \fn{ip_regression.json} in \fn{results/regression/}; full mode).

\begin{table}[ht]
\centering\small
\begin{tabular}{@{}lcc@{}}
\toprule
instance & expected & obtained\\
\midrule
$\slack(\Tseven)$; a pairing for $(\Tseven,\emptyset)$ exists & $1$; yes & $1$; yes\\
$\worst(\Tseven,\emptyset,7)$ & $6$ & $6$\\
$\slack(\Teight)$ & $8$ & $8$\\
$\worst(\Teight,\emptyset,8)$ & $0$ & $0$\\
the same without the constraints \eqref{eq:ip-q}, \eqref{eq:ip-danger} & $1$ & $1$\\
$\worst(\Teight,\emptyset,9)$ & $1$ & $1$\\
$\worst(\Teight,\{h\},8)$ for each of the $16$ busy cells $h\notin\Teight$ & $0$ & $0$ ($16$ times)\\
$\worst(\Teight-(5,5,5)+(1,1,1),\{(5,5,5)\},8)$ & $0$ & $0$\\
antipode stealing after $8$, $9$, $10$ rounds & $6$, $0$, $0$ & $6$, $0$, $0$\\
threat branch: $\Teight\cup\{(0,1,2)\}$, $F=\{(x,1,2):x\ge1\}$, $k=9$ & $0$ & $0$\\
threat branch: $\Tseven\cup\{(0,1,2)\}$, the same $F$, $k=8$ & $2$ & $2$\\
\bottomrule
\end{tabular}
\caption{Regression values for the integer program. The instances are those of \cref{tab:leafvectors}; the
three antipode-stealing values were also confirmed by exact enumeration.}
\label{tab:ip-regression}
\end{table}

In addition, $20$ random sets of eight busy cells (seeded) all fail, with $\worst$ between $2$ and $12$:
a set of eight busy cells is typically far from admitting a pairing strategy, which makes the six orbits of
\cref{ssec:leaftable} remarkable. On seven small instances the program was compared with a brute-force
enumeration of $\worst$ (for example $\worst(\emptyset,\emptyset,0)=80$ and $\worst(\emptyset,\emptyset,1)=81$
on the empty board), with equal values. On all $45$ regression instances the leaf condition is consistent
with the integer program: $\LEAF$ is true on $21$ of them, and these are exactly the instances with
$\worst\le0$.

\subsection{Running times and use}\label{ssec:ip-times}

The running time depends strongly on the number $\theta=k-|F|$ of free stones and on the answer. Instances
that fail are usually fast, because a violating position is found early. Instances that hold require a
proof of optimality. On the leaves of the certificate (decide mode, all answers ``holds'') the median
time per instance was $0.03$\,s and $0.10$\,s for $\theta=0,1$; $1.8$\,s and $2.8$\,s for $\theta=2,3$;
about $12.5$\,s for $\theta=4,5$ (maximum $52$\,s); $117$\,s for $\theta=6$ (maximum $931$\,s); and
$77$\,s for $\theta=7$ (four instances, between $37$ and $133$\,s). In full mode $\worst(\Teight,\emptyset,8)=0$ took $19$\,s, the sixteen instances
$\worst(\Teight,\{h\},8)$ between $26$ and $69$\,s, and
$\worst(\Teight-(5,5,5)+(1,1,1),\{(5,5,5)\},8)$ took $155$\,s. On the leaf table the six orbits for $k=8$
needed between $18$ and $297$\,s each and the eight orbits for $k=9$ that hold between $97$ and $323$\,s. In
short: seconds for small $\theta$, minutes for $\theta\ge6$, and occasionally more; with a time limit of
$300$\,s four of the $1{,}279$ non-degenerate instances of the comparison in \cref{ssec:supp-cross} remained
undecided.

The coverage achieved with the integer program is stated in \cref{ssec:supp-cross}: all $478$ leaves with
$\theta\ge4$ and $830$ sampled leaves with $\theta\le3$; the leaf table for $k=8$ and part of it for $k=9$;
the regression instances; and $1{,}275$ of $1{,}291$ test instances compared with an exact integer decision
procedure. In all cases in which both the integer program and an exact method gave an answer, the answers
agree. We repeat that all ``holds'' answers of the integer program rest on floating-point bounds; they
corroborate the certificate but prove nothing, and \cref{thm:main} does not depend on them.
